\documentclass[12pt,a4paper]{article}
\usepackage[utf8]{inputenc}
\usepackage[T1]{fontenc}
\usepackage[english]{babel}
\usepackage{amsmath}
\usepackage{amssymb}
\usepackage{booktabs}
\usepackage{longtable}
\usepackage{geometry}
\usepackage{hyperref}
\usepackage{xcolor}
\usepackage{listings}
\usepackage{caption}
\usepackage{array}
\usepackage{multirow}

\geometry{margin=2.5cm}

\hypersetup{
    colorlinks=true,
    linkcolor=blue,
    urlcolor=blue,
    citecolor=blue
}

\lstset{
    basicstyle=\ttfamily\small,
    backgroundcolor=\color{gray!10},
    frame=single,
    breaklines=true,
    keywordstyle=\color{blue},
    commentstyle=\color{green!60!black},
    stringstyle=\color{red!70!black},
}

\title{\textbf{Di-Muon Analysis Documentation in Heavy-Ion Collisions}\\
\large FastFrames Configuration --- \texttt{ttmumu.yaml}}
\author{ATLAS Heavy-Ion Analysis}
\date{March 2026}

\begin{document}

\maketitle
\tableofcontents
\newpage

%=============================================================================
\section{Overview}
%=============================================================================

This documentation describes the procedures, selection cuts, and variables
used in the di-muon pair analysis in heavy-ion (Heavy-Ion) collisions
with the ATLAS detector. The analysis uses the \textbf{FastFrames} framework based on
\texttt{RDataFrame} from ROOT for vectorised processing of ntuples.

The main goal is to produce histograms for spectrum fitting with
TRex-Fitter, targeting the study of $Z \to \mu\mu$ production in Pb+Pb collisions.

%=============================================================================
\section{General Configuration}
%=============================================================================

\subsection{Run Parameters}

\begin{table}[h!]
\centering
\caption{General configuration parameters.}
\begin{tabular}{ll}
\toprule
\textbf{Parameter} & \textbf{Value} \\
\midrule
Debug level                    & \texttt{DEBUG} \\
Reconstruction tree            & \texttt{reco} \\
Number of CPUs                 & 4 \\
Automatic systematics          & Disabled \\
Nominal only                   & Yes \\
Region subfolders              & No \\
\bottomrule
\end{tabular}
\end{table}

\subsection{Luminosity and Campaign}

The analysis uses data from the \textbf{mc23d} campaign, with integrated luminosity:

\begin{equation}
    \mathcal{L}_{\text{int}} = 0.00163\ \text{pb}^{-1} = 1630\ \mu\text{b}^{-1}
\end{equation}

\subsection{Event Weights}

The default weight applied to each simulated event is:

\begin{equation}
    w_{\text{event}} = w_{\text{mc}} \times w_{\text{mc,FCal}} \times w_{\text{leptonSF,tight}}
\end{equation}

where:
\begin{itemize}
    \item $w_{\text{mc}}$ (\texttt{weight\_mc\_NOSYS}): Monte Carlo generation weight;
    \item $w_{\text{mc,FCal}}$ (\texttt{weight\_mc\_fcal}): centrality reweighting via FCal;
    \item $w_{\text{leptonSF,tight}}$ (\texttt{weight\_leptonSF\_tight\_NOSYS}): lepton efficiency scale factor (\textit{Tight} working point).
\end{itemize}

\subsection{Samples}

\begin{table}[h!]
\centering
\caption{Samples used in the analysis.}
\begin{tabular}{llll}
\toprule
\textbf{Name} & \textbf{DSID} & \textbf{Campaign} & \textbf{Type} \\
\midrule
\texttt{data}     & 0      & 2024  & Real data \\
\texttt{ttHijing} & 601589 & mc23d & Full Sim (MC) \\
\texttt{zmumu}    & 600646 & mc23d & Full Sim (MC) \\
\bottomrule
\end{tabular}
\end{table}

The cross-section file used is \texttt{PMGxsecDB\_mc23.txt},
which contains the cross-section, $k$-factor, and filter efficiency for each MC sample.

%=============================================================================
\section{Derived Variable Definitions (\textit{Custom Columns})}
%=============================================================================

The variables below are computed at runtime via
\texttt{RDataFrame::Define}, using valid C++ syntax. All momentum quantities
are in \textbf{MeV}, following the ATLAS ntuple convention.

\subsection{Muon Counting}

\begin{equation}
    n_{\mu}^{\text{tight}} = \left| \left\{ \mu_i : p_{T,i} \geq 15000\ \text{MeV}\ \wedge\ \texttt{mu\_select\_tight}_i = \text{true} \right\} \right|
\end{equation}

\noindent Implemented as:
\begin{lstlisting}[language=C++]
mu_pt_NOSYS[mu_pt_NOSYS >= 15000 && mu_select_tight_NOSYS == true].size()
\end{lstlisting}

\subsection{Muon Sorting by $p_T$}

Tight muons are sorted in descending $p_T$ order using
\texttt{ROOT::VecOps::Argsort}:

\begin{lstlisting}[language=C++]
// Sorting index (descending):
ROOT::VecOps::Argsort(-mu_pt_NOSYS[mu_select_tight_NOSYS])

// Sorted variables (example for pT):
ROOT::VecOps::Take(mu_pt_NOSYS[mu_select_tight_NOSYS], mu_sort_idx_NOSYS)
\end{lstlisting}

The following variables are created for pT-sorted muons:

\begin{table}[h!]
\centering
\caption{Derived variables for muons sorted by $p_T$ (momentum in MeV; $\eta$, $\phi$, charge are dimensionless).}
\begin{tabular}{ll}
\toprule
\textbf{Variable} & \textbf{Description} \\
\midrule
\texttt{mu\_pt\_sorted\_NOSYS}         & $p_T$ of sorted Tight muons \\
\texttt{mu\_eta\_sorted\_NOSYS}        & $\eta$ of sorted Tight muons \\
\texttt{mu\_phi\_sorted\_NOSYS}        & $\phi$ of sorted Tight muons \\
\texttt{mu\_e\_sorted\_NOSYS}          & Energy of sorted Tight muons \\
\texttt{mu\_charge\_sorted\_NOSYS}     & Electric charge of sorted Tight muons \\
\texttt{mu\_pt2cone20\_sorted\_NOSYS}  & Track isolation $p_T^{\text{cone20}}$ (MeV), sorted \\
\texttt{mu\_topoetcone20\_sorted\_NOSYS} & Calorimeter isolation $E_T^{\text{topocone20}}$ (MeV), sorted \\
\bottomrule
\end{tabular}
\end{table}

\subsection{Muon Four-Vectors}

The four-vectors of the two highest-$p_T$ muons are constructed using
\texttt{ROOT::Math::PtEtaPhiEVector} ($(p_T, \eta, \phi, E)$ system):

\begin{lstlisting}[language=C++]
// Leading muon (index 0):
ROOT::Math::PtEtaPhiEVector(mu_pt_sorted_NOSYS[0],
                             mu_eta_sorted_NOSYS[0],
                             mu_phi_sorted_NOSYS[0],
                             mu_e_sorted_NOSYS[0])

// Sub-leading muon (index 1):
ROOT::Math::PtEtaPhiEVector(mu_pt_sorted_NOSYS[1],
                             mu_eta_sorted_NOSYS[1],
                             mu_phi_sorted_NOSYS[1],
                             mu_e_sorted_NOSYS[1])
\end{lstlisting}

The invariant mass of the di-muon pair is computed as:

\begin{equation}
    m_{\mu\mu} = \left( P_{\mu_0} + P_{\mu_1} \right)^2, \quad
    \text{via } \texttt{(mu\_tlv0\_NOSYS + mu\_tlv1\_NOSYS).M()}
\end{equation}

%=============================================================================
\section{Selection Regions}
%=============================================================================

The analysis defines four regions, each with a specific event selection:

%-----------------------------------------------------------------------------
\subsection{Region \texttt{all\_tight\_muon} --- Diagnostic}
%-----------------------------------------------------------------------------

General control region used for diagnostics. Requires only the event
selection flag and exactly two Tight muons.

\begin{table}[h!]
\centering
\caption{Cuts for the \texttt{all\_tight\_muon} region.}
\begin{tabular}{lll}
\toprule
\textbf{Variable} & \textbf{Cut} & \textbf{Description} \\
\midrule
\texttt{pass\_mumu\_NOSYS}                        & $= 1$     & Di-muon event selection flag \\
$\sum_i \texttt{mu\_select\_tight}_i$             & $= 2$     & Exactly 2 Tight muons \\
\bottomrule
\end{tabular}
\end{table}

%-----------------------------------------------------------------------------
\subsection{Region \texttt{sumET} --- Centrality Distribution}
%-----------------------------------------------------------------------------

No selection cut (accepts all events). Used to monitor the
FCal transverse energy distribution, which characterises the collision
centrality.

\begin{equation}
    \text{Selection: no cut (}\texttt{selection: "1"}\text{)}
\end{equation}

%-----------------------------------------------------------------------------
\subsection{Region \texttt{TREXMU} --- Main Di-Muon Selection}
%-----------------------------------------------------------------------------

Main analysis region with a full di-muon selection for TRex-Fitter fitting.

\begin{table}[h!]
\centering
\caption{Cuts for the \texttt{TREXMU} region (full di-muon selection).}
\begin{tabular}{lll}
\toprule
\textbf{Variable} & \textbf{Cut} & \textbf{Description} \\
\midrule
\texttt{pass\_mumu\_NOSYS}                          & $= 1$              & Di-muon event selection flag \\
$n_\mu^{\text{tight}}$                              & $= 2$              & Exactly 2 Tight muons with $p_T > 15$~GeV \\
\texttt{fcal\_sum\_et}                              & $> 0.063208$~TeV   & FCal centrality cut ($< 80\%$) \\
$|\eta_{\mu_0}|$                                    & $< 2.47$           & Pseudorapidity acceptance (leading) \\
$|\eta_{\mu_1}|$                                    & $< 2.47$           & Pseudorapidity acceptance (sub-leading) \\
$q_{\mu_0} \cdot q_{\mu_1}$                         & $< 0$              & Opposite-charge pair \\
$p_{T,\mu_0}^{\text{cone20}} / p_{T,\mu_0}$        & $< 0.32$           & Track isolation (leading) \\
$p_{T,\mu_1}^{\text{cone20}} / p_{T,\mu_1}$        & $< 0.32$           & Track isolation (sub-leading) \\
$E_{T,\mu_0}^{\text{topocone20}} / p_{T,\mu_0}$    & $< 0.5$            & Calorimetric isolation (leading) \\
$E_{T,\mu_1}^{\text{topocone20}} / p_{T,\mu_1}$    & $< 0.5$            & Calorimetric isolation (sub-leading) \\
$\sum_i \texttt{el\_select\_tight}_i$               & $= 0$              & Tight electron veto \\
\bottomrule
\end{tabular}
\end{table}

\subsubsection*{Notes on Isolation}

Isolation cuts are applied as dimensionless ratios, with both numerator
and denominator in MeV. For the leading muon ($\mu_0$):

\begin{align}
    \frac{p_T^{\text{cone20}}}{p_T} &< 0.32 \quad \text{(track isolation)} \\[6pt]
    \frac{E_T^{\text{topocone20}}}{p_T} &< 0.5  \quad \text{(calorimetric isolation)}
\end{align}

\subsubsection*{Note on Centrality}

The cut $E_T^{\text{FCal}} > 0.063208$~TeV corresponds to collisions with
centrality $< 80\%$, selecting only collisions with significant geometric
overlap between the nuclei.

%-----------------------------------------------------------------------------
\subsection{Region \texttt{TREXET\_2MU\_CENTRAL\_0\_80} --- Central Di-Muon}
%-----------------------------------------------------------------------------

Applies exactly the same cuts as the \texttt{TREXMU} region, but the
histograms produced are the FCal $\Sigma E_T$ variables instead of muon
variables. Used to study the centrality distribution within the di-muon selection.

%=============================================================================
\section{Produced Histograms}
%=============================================================================

\subsection{Common Variables (regions \texttt{all\_tight\_muon} and \texttt{TREXMU})}

\begin{longtable}{llccc}
\caption{Histogrammed variables in the muon regions.} \\
\toprule
\textbf{Name} & \textbf{Definition} & \textbf{Min} & \textbf{Max} & \textbf{Bins} \\
\midrule
\endfirsthead
\toprule
\textbf{Name} & \textbf{Definition} & \textbf{Min} & \textbf{Max} & \textbf{Bins} \\
\midrule
\endhead
\texttt{mu\_pt}           & \texttt{mu\_pt\_NOSYS} [MeV]                         & 0      & 200\,000 & 20  \\
\texttt{mu\_eta}          & \texttt{mu\_eta}                                      & $-3.0$ & $3.0$    & 60  \\
\texttt{n\_muons}         & \texttt{nMuons\_NOSYS}                                & $-0.5$ & $7.5$    & 8   \\
\texttt{mu\_charge}       & \texttt{mu\_charge}                                   & $-2.0$ & $2.0$    & 4   \\
\texttt{m\_mumu}          & $(P_{\mu_0}+P_{\mu_1}).\texttt{M()}$ [MeV]           & 50\,000 & 120\,000 & 100 \\
\texttt{ratio\_pt2cone20} & $p_T^{\text{cone20}}[0] / p_T[0]$                    & 0      & 1        & 50  \\
\texttt{ratio\_topoetcone20} & $E_T^{\text{topocone20}}[0] / p_T[0]$             & 0      & 1        & 50  \\
\texttt{mu\_pt2cone20}    & \texttt{mu\_pt2cone20\_sorted\_NOSYS[0]} [MeV]        & 0      & 50       & 50  \\
\bottomrule
\end{longtable}

\noindent\textbf{Note:} The $m_{\mu\mu}$ histogram range ($[50\,000,\ 120\,000]$ MeV)
corresponds to the $Z$ boson mass window, approximately $[50,\ 120]$ GeV.

\subsection{Centrality Variables (regions \texttt{sumET} and \texttt{TREXET\_2MU\_CENTRAL\_0\_80})}

\begin{table}[h!]
\centering
\caption{Histogrammed variables in the transverse energy regions.}
\begin{tabular}{llccc}
\toprule
\textbf{Name} & \textbf{Definition} & \textbf{Min} & \textbf{Max} & \textbf{Bins} \\
\midrule
\texttt{Fcal\_sum\_ET} & \texttt{fcal\_sum\_et} [TeV] & 0 & 6 & 60 \\
\bottomrule
\end{tabular}
\end{table}

%=============================================================================
\section{Output Ntuples}
%=============================================================================

In addition to histograms, filtered ntuples are produced for the
\texttt{TREXMU} and \texttt{TREXET\_2MU\_CENTRAL\_0\_80} regions. The saved branches are:

\begin{itemize}
    \item \texttt{.*\_pt\_NOSYS} --- all $p_T$ branches (muons, electrons, jets);
    \item \texttt{mu\_eta} --- muon pseudorapidity;
    \item \texttt{mu\_phi} --- muon azimuthal angle.
\end{itemize}

%=============================================================================
\section{Selection Flow}
%=============================================================================

The following steps summarise the event selection flow:

\begin{enumerate}
    \item \textbf{Pre-selection:} \texttt{pass\_mumu\_NOSYS == 1}
    \item \textbf{Muon quality:} select muons passing the \textit{Tight} working point
          with $p_T > 15$~GeV
    \item \textbf{Sorting:} rank Tight muons by decreasing $p_T$
    \item \textbf{Centrality:} $E_T^{\text{FCal}} > 0.063208$~TeV (centrality $< 80\%$)
    \item \textbf{Multiplicity:} exactly $N_\mu^{\text{tight}} = 2$
    \item \textbf{Acceptance:} $|\eta_{\mu_{0,1}}| < 2.47$
    \item \textbf{Opposite charge:} $q_{\mu_0} \cdot q_{\mu_1} < 0$
    \item \textbf{Track isolation:} $p_T^{\text{cone20}} / p_T < 0.32$ (both muons)
    \item \textbf{Calorimetric isolation:} $E_T^{\text{topocone20}} / p_T < 0.5$ (both)
    \item \textbf{Electron veto:} $N_e^{\text{tight}} = 0$
\end{enumerate}

%=============================================================================
\section{Technical Notes}
%=============================================================================

\begin{itemize}
    \item The trigger \texttt{HLT\_mu4\_L1MU3V} was identified as the only active trigger
          in HI data. In MC, the trigger branch often does not exist;
          therefore, the trigger cut is not applied in the selection (MC is $\sim$100\%
          efficient for muons with $p_T > 15$~GeV).
    \item All momenta in ATLAS HI ntuples are in \textbf{MeV}.
    \item The four-vector construction uses \texttt{ROOT::Math::PtEtaPhiEVector}
          with MeV inputs; the resulting invariant mass $m_{\mu\mu}$ is therefore
          also in MeV.
    \item The \texttt{nominal\_only: True} option disables systematic variation
          processing, running only the nominal configuration.
    \item The framework does not use \texttt{ROOT::RVec} directly
          (\texttt{use\_rvec: False}); vectors are processed via
          \texttt{ROOT::VecOps}.
\end{itemize}

\end{document}
